\section{A flux-only count for globally propagating vacuum fields}
\label{cb:section}

This section gives a second kind of restricted channel theorem: the entire
planar fields must arise from propagating vacuum Herglotz waves. Their
amplitudes may be arbitrarily large relative to the chosen input norm.
Absolute bypass then supplies a finite dimension bound without an electric
or magnetic field-resource bound. The conclusion is not an area law and
does not apply to general dielectric fields or local evanescent solutions.
It also identifies an obstruction to completing a remote radial-control
construction by sending controls to infinity around a fixed regular-wave
target.

\subsection{An explicit scalar observation estimate}

Use the Fourier convention
$\widehat f(\xi)=\int f(x)e^{-i\xi\cdot x}\,dx$.
For $\kappa,L>0$, define the explicit numbers
\begin{equation}
 n=2+\lceil4e\kappa L\rceil,\qquad
 C_{\kappa,L}=2\left[1+2n(16e)^{n-1}\right].
 \label{cb:observationconstant}
\end{equation}
Here $e$ in the numerical constants is the base of the natural logarithm.

\begin{lemma}[PROVED: observation outside an interval]
\label{cb:onedimensional}
If $f\in L^1(\mathbb R)$ and
$\operatorname{supp}\widehat f\subset[-\kappa,\kappa]$, then, for
$I=(-L/2,L/2)$,
\begin{equation}
 \|f\|_{L^1(\mathbb R)}
 \leq C_{\kappa,L}\|f\|_{L^1(\mathbb R\setminus I)}.
 \label{cb:onedimensionalbound}
\end{equation}
\end{lemma}
\begin{proof}
Fourier inversion applies because $\widehat f$ is bounded and compactly
supported. In particular $f$ has a smooth representative, and
\begin{equation}
 \|f^{(j)}\|_\infty
 \leq\frac{\kappa^{j+1}}{\pi(j+1)}\|f\|_1
 \quad(j\geq0).
 \label{cb:onedimensionalderivative}
\end{equation}
Write $F=\|f\|_1$ and $G=\|f\|_{L^1(\mathbb R\setminus I)}$.
In the adjacent interval $[L/2,3L/2]$, take the $n$ subintervals
\begin{equation}
 J_j=\left[L/2+jL/n,\ L/2+(j+1/2)L/n\right],
 \qquad 0\leq j<n.
 \label{cb:nodeslots}
\end{equation}
Each has length $h=L/(2n)$. Averaging gives a point $x_j\in J_j$ with
$|f(x_j)|\leq G/h=2nG/L$. For $j>i$,
$x_j-x_i\geq(j-i)h$. Let $P$ be the polynomial of degree at most
$n-1$ interpolating these values. For $x\in I$, all distances
$|x-x_j|$ are at most $2L$, so its Lagrange formula gives
\begin{align}
 \sum_{j=0}^{n-1}
  \left|\prod_{i\ne j}\frac{x-x_i}{x_j-x_i}\right|
 &\leq\frac{(4n)^{n-1}2^{n-1}}{(n-1)!}
 \leq(16e)^{n-1},\notag\\
 |P(x)|&\leq\frac{2n}{L}(16e)^{n-1}G.
 \label{cb:lagrangebound}
\end{align}
The last inequality uses $m!\geq(m/e)^m$ with $m=n-1\geq1$.
The divided-difference integral remainder, obtained by repeated integration
of the fundamental theorem of calculus, applies also to complex-valued $f$:
\begin{equation}
 |f(x)-P(x)|\leq\frac{(2L)^n}{n!}\|f^{(n)}\|_\infty.
 \label{cb:interpolationremainder}
\end{equation}
It uses derivative values only on the real convex hull of the nodes and $x$;
no complex mean-value theorem is being assumed.
Integrating over $I$ and using \eqref{cb:onedimensionalderivative},
\begin{equation}
 \|f\|_{L^1(I)}
 \leq2n(16e)^{n-1}G+
    \frac{\kappa L}{\pi(n+1)}\frac{(2\kappa L)^n}{n!}F.
 \label{cb:absorption}
\end{equation}
Our choice $n\geq4e\kappa L$ implies
$(2\kappa L)^n/n!\leq2^{-n}$. Hence the coefficient of $F$ is
at most $2^{-n}/(4e\pi)<1/2$. Add $G$ to both sides of
\eqref{cb:absorption} and absorb this term to obtain
\eqref{cb:onedimensionalbound}.
\end{proof}

\begin{corollary}[PROVED: planar observation outside the fixed square]
\label{cb:planarobservation}
Let $D=(-L/2,L/2)^2$. If $f\in L^1(\mathbb R^2)$ and
$\operatorname{supp}\widehat f\subset\{\xi:|\xi|\leq\kappa\}$, then
\begin{equation}
 \|f\|_{L^1(\mathbb R^2)}
 \leq C_{\kappa,L}\|f\|_{L^1(\mathbb R^2\setminus D)},\qquad
 \|\nabla f\|_\infty
 \leq\frac{\kappa^3}{6\pi}\|f\|_1.
 \label{cb:planarestimates}
\end{equation}
\end{corollary}
\begin{proof}
Almost every horizontal slice is in $L^1(\mathbb R)$ and has one-dimensional
Fourier support in $[-\kappa,\kappa]$. To justify the support statement,
test its partial Fourier transform against a smooth function supported
outside that interval. Fourier transformation in the second variable gives
zero by the global support assumption; injectivity and Fubini give zero for
almost every slice. A countable dense collection of tests suffices.
Apply Lemma~\ref{cb:onedimensional} on these slices and integrate in the
second variable. The observed strip $\{|x|\geq L/2\}$ is contained in
$\mathbb R^2\setminus D$, proving the first assertion. For the second,
Fourier inversion and $|\widehat f|\leq\|f\|_1$ give
\begin{equation}
 |\nabla f(x)|\leq\frac{\|f\|_1}{(2\pi)^2}
       \int_{|\xi|\leq\kappa}|\xi|\,d\xi
   =\frac{\kappa^3}{6\pi}\|f\|_1.
 \label{cb:planarderivative}
\end{equation}
\end{proof}

\subsection{Finite positive trace from actual pointwise flux}

\begin{theorem}[PROVED: channel finiteness from propagating bandwidth and bypass]
\label{cb:count}
Let $W$ be a complex linear input space equipped with a Hilbert norm. Suppose
its electric and magnetic fields depend linearly on the input, and are
entire regular vacuum Maxwell Herglotz fields at real wavenumber $k>0$.
There is no assumed bound on their Herglotz density norms in terms of
$\|a\|$. Suppose every coherent $a\in W$ satisfies
\begin{equation}
 P_D(a)=\int_Ds_a\geq p\|a\|^2,\qquad
 \int_{\mathbb R^2\setminus D}|s_a|\leq\delta\|a\|^2,
 \quad
 s_a=\tfrac12\operatorname{Re}(E_a\times\overline{H_a})\cdot\widehat z,
 \label{cb:fluxhypotheses}
\end{equation}
where $D=(-L/2,L/2)^2$, $p>0$, and $\delta\geq0$. Set $\kappa=2k$ and
$A=L^2$. Then
\begin{equation}
 \dim W\leq\frac{A\kappa^2 C_{\kappa,L}\delta}{\pi p}
          =\frac{4Ak^2 C_{2k,L}\delta}{\pi p}.
 \label{cb:dimension}
\end{equation}
If $\delta=0$, in fact $W=\{0\}$. For a counted physical subspace one
may take $p=\eta\tau$; no condition $\delta<\eta\tau$ is needed here.
\end{theorem}
\begin{proof}
On the plane each electric and magnetic component is the Fourier transform
of a finite complex measure supported in the transverse disk of radius $k$.
Indeed, push its spherical Herglotz density forward under
$d\mapsto kd_\perp$; the measure is finite because $L^2(S^2)\subset L^1(S^2)$.
Products of a component and a conjugate component correspond to convolution
with a reflected conjugate measure. Thus
\begin{equation}
 \operatorname{supp}\widehat{s_a}\subset\{|\xi|\leq2k\}.
 \label{cb:fluxbandwidth}
\end{equation}
The fields are continuous, hence $s_a$ is integrable on the bounded square.
Its assumed absolute bypass implies global $L^1$ integrability. In particular,
global flux integrability is not being inferred from an arbitrary
$L^2$ Herglotz density alone.

Corollary~\ref{cb:planarobservation} now gives
\begin{equation}
 \|s_a\|_1\leq C_{\kappa,L}\delta\|a\|^2,\qquad
 |s_a(x)|\leq\frac{\kappa^2 C_{\kappa,L}\delta}{4\pi}\|a\|^2.
 \label{cb:fluxcontrol}
\end{equation}
The pointwise bound follows from Fourier inversion, since the frequency
disk has area $\pi\kappa^2$ and $|\widehat{s_a}|\leq\|s_a\|_1$.
This controls the absolute flux \emph{inside} $D$ as well as outside it;
the signed passage alone would not do so.

Restrict to an arbitrary finite-dimensional input subspace $V\subset W$
of dimension $N$. At each point $x$, write its Hermitian normal-flux form
as $s_a(x)=\langle a,Q(x)a\rangle$. This form factors through the linear map
\begin{equation}
 a\longmapsto\bigl(E_{a,x}(x),E_{a,y}(x),H_{a,x}(x),H_{a,y}(x)\bigr)
       \in\mathbb C^4.
 \label{cb:pointmap}
\end{equation}
Consequently $\operatorname{rank}Q(x)\leq4$, and
\eqref{cb:fluxcontrol} bounds its operator norm by
$\kappa^2 C_{\kappa,L}\delta/(4\pi)$. Thus
\begin{equation}
 \operatorname{Tr}Q(x)\leq\operatorname{Tr}(Q(x))_+
       \leq4\|Q(x)\|
       \leq\frac{\kappa^2 C_{\kappa,L}\delta}{\pi}.
 \label{cb:pointtrace}
\end{equation}
For an orthonormal basis $a_1,\ldots,a_N$ of $V$, the passage hypothesis gives
\begin{equation}
 pN\leq\sum_{j=1}^NP_D(a_j)
    =\int_D\operatorname{Tr}Q(x)\,dx
    \leq\frac{A\kappa^2 C_{\kappa,L}\delta}{\pi}.
 \label{cb:integratedtrace}
\end{equation}
Only a finite sum is interchanged with integration. This proves
\eqref{cb:dimension}, including infinite-dimensional $W$ by restriction
to finite-dimensional subspaces.
If $\delta=0$, \eqref{cb:fluxcontrol} gives $s_a=0$, and positive passage
excludes every nonzero input.
\end{proof}

\begin{remark}[Scope of the explicit estimate]
\label{cb:scope}
This is a flux-resource estimate, not a field-energy estimate. It permits
unbounded electric and magnetic field norms and arbitrary coherent
interference. The explicit observation constant grows exponentially in
$kL$, and the resulting bound is not the proposed area scaling.
Neither local vacuum Maxwell equations nor the existence of an entire
extension alone implies \eqref{cb:fluxbandwidth}: entire solutions with
evanescent transverse content can grow in the normal direction and fall
outside the propagating Herglotz class. The earlier local packet examples
therefore do not contradict this theorem. General fields on a plane
intersecting a dielectric also fall outside its hypotheses.
\end{remark}

\subsection{The limitation of moving controls to infinity}

\begin{lemma}[PROVED: uniform tails in a fixed regular-wave band]
\label{cb:finitetail}
Fix a positive integer $J$, and let $\mathcal F_J$ be the real linear space
of planar Hermitian cross-flux densities of regular vacuum vector spherical
waves with angular orders $1\leq l\leq J$, both polarizations and all
azimuthal indices. Then $\mathcal F_J$ is finite dimensional and contained
in $L^1(\mathbb R^2)$. For every $\epsilon>0$ there is a finite radius
$R_{J,\epsilon}$ such that
\begin{equation}
 \int_{|x|>R_{J,\epsilon}}|f(x)|\,dx
       \leq\epsilon\|f\|_1\qquad(f\in\mathcal F_J).
 \label{cb:uniformtail}
\end{equation}
\end{lemma}
\begin{proof}
There are finitely many vector waves, so their Hermitian cross-fluxes span
a finite-dimensional function space. They are smooth near the origin.
Their finite spherical Hankel formulas give, uniformly in direction,
\begin{equation}
 \begin{split}
 E(r\widehat x)&=r^{-1}\bigl(e^{ikr}A_+(\widehat x)
                         +e^{-ikr}A_-(\widehat x)\bigr)+O(r^{-2}),\\
 H(r\widehat x)&=(Z_0r)^{-1}\widehat x\times
                 \bigl(e^{ikr}A_+(\widehat x)
                         -e^{-ikr}A_-(\widehat x)\bigr)+O(r^{-2}),
 \end{split}
 \label{cb:regularasymptotic}
\end{equation}
with $A_\pm\cdot\widehat x=0$. For any pair of such waves the leading
cross-Poynting vector is radial, by
$A\times(\widehat x\times\overline B)
=\widehat x(A\cdot\overline B)$ for transverse $A,B$.
On the measurement plane its normal component therefore vanishes.
Each cross-flux is $O(r^{-3})$, which is integrable in the two-dimensional
plane. For every fixed $f\in\mathcal F_J$ its $L^1$ tail tends to zero.
The unit sphere of the actual function space $\mathcal F_J$, with its
$L^1$ norm, is compact. The tail functionals are all Lipschitz with
constant one in that norm. A finite net of the unit sphere consequently
makes this convergence uniform, proving \eqref{cb:uniformtail}.
\end{proof}

\begin{corollary}[PROVED: finite-band control beyond a sufficiently large vacuum core]
\label{cb:finitecore}
Fix $J$ and choose
$R_J\geq L/\sqrt2$ so that \eqref{cb:uniformtail} holds with
$\epsilon=1/(2C_{2k,L})$. Consider any device whose material is outside
the ball $B_{R_J}$ and whose fields on its counted input space, within that
ball, are linear combinations of regular vacuum vector waves of orders
$l\leq J$. If the actual fields satisfy \eqref{cb:fluxhypotheses}, then
\begin{equation}
 \dim W\leq\frac{2A(2k)^2C_{2k,L}\delta}{\pi p}.
 \label{cb:finitecorecount}
\end{equation}
There is no bound on the coefficients or gains in this finite band. The
radius threshold $R_J$ may depend on the band, while the displayed count
does not.
\end{corollary}
\begin{proof}
Extend each regular-wave combination to an entire vacuum field, and let
$f_a\in\mathcal F_J$ be its plane flux. It agrees with the actual flux in
the radius-$R_J$ disk. Corollary~\ref{cb:planarobservation} and the uniform
tail estimate imply
\begin{align}
 \|f_a\|_1
 &\leq C_{2k,L}\int_{\mathbb R^2\setminus D}|f_a|\notag\\
 &\leq C_{2k,L}\int_{\{|x|\leq R_J\}\setminus D}|s_a|
           +\tfrac12\|f_a\|_1
 \leq C_{2k,L}\delta\|a\|^2+\tfrac12\|f_a\|_1.
 \label{cb:tailabsorption}
\end{align}
Hence $\|f_a\|_1\leq2C_{2k,L}\delta\|a\|^2$.
The pointwise Fourier and rank-four trace argument in
\eqref{cb:pointtrace}--\eqref{cb:integratedtrace} applies with this doubled
constant, while passage is unchanged on $D$. This proves
\eqref{cb:finitecorecount}.
\end{proof}

\begin{corollary}[PROVED under explicit convergence: fixed-target remote-control obstruction]
\label{cb:remotelimit}
Let a sequence of physical devices be vacuum on balls with radii tending
to infinity. Identify a fixed finite-dimensional normed input space $W$
for the sequence. Assume that, for every $a\in W$, its electric and magnetic
fields converge uniformly on compact subsets of the measurement plane
to an entire propagating vacuum Herglotz family $(E_a^0,H_a^0)$.
If all devices satisfy \eqref{cb:fluxhypotheses} with the same $p,\delta$, then
their fixed input dimension obeys \eqref{cb:dimension}.
\end{corollary}
\begin{proof}
Compact uniform convergence gives convergence of the flux on $D$, so
$\int_Ds_a^0\geq p\|a\|^2$. Pointwise convergence on the complement and
Fatou's lemma give
$\int_{\mathbb R^2\setminus D}|s_a^0|\leq\delta\|a\|^2$.
The limit satisfies all hypotheses of Theorem~\ref{cb:count}.
\end{proof}

\begin{remark}[OPEN: varying targets and nearby material]
\label{cb:remotescope}
The convergence and propagating-limit hypotheses in
Corollary~\ref{cb:remotelimit} are essential parts of its statement.
Finite-band transfer control, by itself, supplies no global bypass estimate.
Allowing the angular band to grow with the remote-control radius, together
with changing coefficients, can produce local superoscillatory limits with
large vacuum tails beyond that radius. No uniform compactness or global
Fourier-support claim for such sequences is made here. Changing coefficients
alone in a fixed band does not evade Corollary~\ref{cb:finitecore} once its
radius threshold is reached. Likewise this argument does not constrain
devices whose material remains near the section. It proves an obstruction
to a fixed-target remote-control completion, while the unrestricted
ordinary-dielectric channel target remains open.
\end{remark}
